% conv2d_matrix.tex — full 9x25 weight matrix with 25x1 input and 9x1 output
% kernel values (row-major, bottom-up as in conv2d.tex):
% row0(bot): 1  0 -2
% row1(mid): 3  3 -1
% row2(top): 0  2  0
% kernel flattened row-major (top-to-bottom): 0,2,0,3,3,-1,1,0,-2
% for output position (i,j), active input cols in 25-element flat vector:
% flat index = row*5 + col
% output row k (k=0..8): i=k/3, j=k%3
% active flat indices: (i+0)*5+(j+0), (i+0)*5+(j+1), (i+0)*5+(j+2),
%                      (i+1)*5+(j+0), (i+1)*5+(j+1), (i+1)*5+(j+2),
%                      (i+2)*5+(j+0), (i+2)*5+(j+1), (i+2)*5+(j+2)

\begin{tikzpicture}[scale=0.75]

\def\cw{0.30}  % cell width
\def\ch{0.30}  % cell height
\def\gap{0.15} % gap between blocks

% kernel values in flat order (top-left to bottom-right of filter)
% kv[0..8]: positions in filter scanned row by row
% filter from conv2d.tex (top row first):
% top:    0  2  0
% mid:    3  3 -1
% bot:    1  0 -2
\def\kva{0}
\def\kvb{2}
\def\kvc{0}
\def\kvd{3}
\def\kve{3}
\def\kvf{-1}
\def\kvg{1}
\def\kvh{0}
\def\kvi{-2}

% active flat input indices for each of 9 output positions
% output k: i=floor(k/3), j=mod(k,3)
% active: (i)*5+j, i*5+j+1, i*5+j+2, (i+1)*5+j, ..., (i+2)*5+j+2

% ─────────────────────────────────────────────────────
% WEIGHT MATRIX  9 rows x 25 cols
% ─────────────────────────────────────────────────────
\foreach \outrow in {0,...,8} {
    % compute patch start: i=floor(outrow/3), j=mod(outrow,3)
    \pgfmathtruncatemacro\pi{int(\outrow/3)}
    \pgfmathtruncatemacro\pj{int(mod(\outrow,3))}
    % 9 active flat indices
    \pgfmathtruncatemacro\aa{int(\pi*5+\pj)}
    \pgfmathtruncatemacro\ab{int(\pi*5+\pj+1)}
    \pgfmathtruncatemacro\ac{int(\pi*5+\pj+2)}
    \pgfmathtruncatemacro\ba{int((\pi+1)*5+\pj)}
    \pgfmathtruncatemacro\bb{int((\pi+1)*5+\pj+1)}
    \pgfmathtruncatemacro\bc{int((\pi+1)*5+\pj+2)}
    \pgfmathtruncatemacro\ca{int((\pi+2)*5+\pj)}
    \pgfmathtruncatemacro\cb{int((\pi+2)*5+\pj+1)}
    \pgfmathtruncatemacro\cc{int((\pi+2)*5+\pj+2)}
    \foreach \incol in {0,...,24} {
        % check if active
        \pgfmathtruncatemacro\act{ifthenelse(
            \incol==\aa || \incol==\ab || \incol==\ac ||
            \incol==\ba || \incol==\bb || \incol==\bc ||
            \incol==\ca || \incol==\cb || \incol==\cc, 1, 0)}
        \ifnum\act=1
            \draw[thwsorange, fill=thwsorange!35, thick]
                (\incol*\cw, -\outrow*\ch)
                rectangle (\incol*\cw+\cw-0.02, -\outrow*\ch+\ch-0.02);
            % kernel value: which position in 3x3 kernel?
            \pgfmathtruncatemacro\kidx{ifthenelse(\incol==\aa,0,
                ifthenelse(\incol==\ab,1,
                ifthenelse(\incol==\ac,2,
                ifthenelse(\incol==\ba,3,
                ifthenelse(\incol==\bb,4,
                ifthenelse(\incol==\bc,5,
                ifthenelse(\incol==\ca,6,
                ifthenelse(\incol==\cb,7,8))))))))}
            \ifnum\kidx=0 \def\kval{\kva} \fi
            \ifnum\kidx=1 \def\kval{\kvb} \fi
            \ifnum\kidx=2 \def\kval{\kvc} \fi
            \ifnum\kidx=3 \def\kval{\kvd} \fi
            \ifnum\kidx=4 \def\kval{\kve} \fi
            \ifnum\kidx=5 \def\kval{\kvf} \fi
            \ifnum\kidx=6 \def\kval{\kvg} \fi
            \ifnum\kidx=7 \def\kval{\kvh} \fi
            \ifnum\kidx=8 \def\kval{\kvi} \fi
            \node[font=\tiny] at
                (\incol*\cw+\cw/2-0.01, -\outrow*\ch+\ch/2-0.01) {\kval};
        \else
            \draw[thwsgrey!25, fill=thwsgrey!8]
                (\incol*\cw, -\outrow*\ch)
                rectangle (\incol*\cw+\cw-0.02, -\outrow*\ch+\ch-0.02);
            \node[font=\tiny, thwsgrey!60] at
                (\incol*\cw+\cw/2-0.01, -\outrow*\ch+\ch/2-0.01) {0};
        \fi
    }
}
\node[font=\scriptsize, thwsorange] at (25*\cw/2, 0.45) {$\mW$ \quad $9\times25$};

% ─────────────────────────────────────────────────────
% INPUT VECTOR  25x1  (to the right of matrix)
% ─────────────────────────────────────────────────────
\pgfmathsetmacro\xinput{25*\cw + \gap}
% patch for output 0: active flat indices 0,1,2,5,6,7,10,11,12
\foreach \i in {0,...,24} {
    \pgfmathtruncatemacro\act{ifthenelse(
        \i==0||\i==1||\i==2||
        \i==5||\i==6||\i==7||
        \i==10||\i==11||\i==12, 1, 0)}
    \ifnum\act=1
        \draw[thwspetrol, fill=thwspetrol!35, thick]
            (\xinput, -\i*\cw) rectangle (\xinput+\cw-0.02, -\i*\cw+\cw-0.02);
    \else
        \draw[thwspetrol!25, fill=thwspetrol!8]
            (\xinput, -\i*\cw) rectangle (\xinput+\cw-0.02, -\i*\cw+\cw-0.02);
    \fi
}
\node[font=\scriptsize, thwspetrol] at (\xinput+\cw/2, 0.45) {$\vx$};
\node[font=\scriptsize, thwspetrol] at (\xinput+\cw/2, -25*\cw-0.2) {$25\times1$};

% multiplication sign
\node[font=\normalsize] at (\xinput-\gap/2, -4*\cw) {$\times$};

% ─────────────────────────────────────────────────────
% OUTPUT VECTOR  9x1  (to the left of matrix)
% ─────────────────────────────────────────────────────
\pgfmathsetmacro\xoutput{-\cw - \gap}
\foreach \i in {0,...,8} {
    \ifnum\i=0
        \draw[carnelian, fill=carnelian!40, thick]
            (\xoutput, -\i*\ch) rectangle (\xoutput+\cw-0.02, -\i*\ch+\ch-0.02);
    \else
        \draw[carnelian!35, fill=carnelian!10]
            (\xoutput, -\i*\ch) rectangle (\xoutput+\cw-0.02, -\i*\ch+\ch-0.02);
    \fi
}
\node[font=\scriptsize, carnelian] at (\xoutput+\cw/2, 0.45) {$\vh$};
\node[font=\scriptsize, carnelian] at (\xoutput+\cw/2, -9*\ch-0.2) {$9\times1$};

% equals sign
\node[font=\normalsize] at (\xoutput+\cw+\gap/2, -4*\cw) {$=$};

% ─────────────────────────────────────────────────────
% ANNOTATIONS
% ─────────────────────────────────────────────────────
\draw[decorate, decoration={brace,amplitude=4pt}, thwsorange]
    (0, 0.35) -- (25*\cw, 0.35);
\node[font=\scriptsize, thwsorange, align=center] at (25*\cw/2, 0.75)
    {same kernel values, shifted --- \textbf{tied weights}};

\draw[decorate, decoration={brace,amplitude=4pt,mirror}, thwsgrey]
    (0, -9*\ch-0.05) -- (25*\cw, -9*\ch-0.05);
\node[font=\scriptsize, thwsgrey] at (25*\cw/2, -9*\ch-0.45)
    {most entries zero --- \textbf{sparse}};

\end{tikzpicture}
